% =====================================================================
%  Tutorial: subir ServiceDesk UTA y ponerlo en marcha en otro servidor
%  Compilar:  latexmk -pdf tutorial_despliegue.tex
%
%  No contiene contraseñas, secretos, direcciones de servidores reales ni
%  datos personales. Los valores entre <...> los sustituye cada instalación.
% =====================================================================
\documentclass[11pt,a4paper,oneside]{report}

\usepackage[utf8]{inputenc}
\usepackage[T1]{fontenc}
\usepackage[spanish,es-noshorthands,es-tabla]{babel}
\usepackage{lmodern}
\usepackage{microtype}
\usepackage[a4paper,top=2.5cm,bottom=2.5cm,left=2.6cm,right=2.4cm,headheight=14pt]{geometry}
\usepackage{fancyhdr}
\usepackage{titlesec}
\usepackage{enumitem}
\setlist{itemsep=2pt,topsep=4pt}
\usepackage[table]{xcolor}
\usepackage{booktabs}
\usepackage{tabularx}
\usepackage{array}
\usepackage{listings}
\usepackage[most]{tcolorbox}
\usepackage{tikz}
\usetikzlibrary{positioning,arrows.meta}
\usepackage{amssymb}
\usepackage{pifont}
\usepackage{xurl}
\usepackage[hidelinks]{hyperref}
\usepackage{bookmark}

\definecolor{utaRojo}{HTML}{8B1A1A}
\definecolor{gris}{HTML}{5A5A5A}
\definecolor{grisClaro}{HTML}{F3F3F3}
\definecolor{grisLinea}{HTML}{CFCFCF}
\definecolor{azulNota}{HTML}{1F4E79}
\definecolor{ambar}{HTML}{8A5A00}
\definecolor{verde}{HTML}{2E6B30}

\hypersetup{colorlinks=true,linkcolor=utaRojo,urlcolor=azulNota,
  pdftitle={Tutorial de despliegue — ServiceDesk UTA},pdfauthor={Equipo ServiceDesk UTA}}

\titleformat{\chapter}[hang]{\normalfont\huge\bfseries\color{utaRojo}}{\thechapter}{0.6em}{}
\titlespacing*{\chapter}{0pt}{-10pt}{16pt}
\titleformat{\section}{\normalfont\Large\bfseries\color{utaRojo}}{\thesection}{0.6em}{}
\titleformat{\subsection}{\normalfont\large\bfseries}{\thesubsection}{0.6em}{}

\pagestyle{fancy}
\fancyhf{}
\fancyhead[L]{\small\color{gris}ServiceDesk UTA · Tutorial de despliegue}
\fancyhead[R]{\small\color{gris}\nouppercase{\leftmark}}
\fancyfoot[C]{\small\thepage}
\renewcommand{\headrulewidth}{0.4pt}
\fancypagestyle{plain}{\fancyhf{}\fancyfoot[C]{\small\thepage}\renewcommand{\headrulewidth}{0pt}}

\lstdefinestyle{shell}{
  basicstyle=\ttfamily\small, backgroundcolor=\color{grisClaro},
  frame=single, rulecolor=\color{grisLinea}, framesep=4pt, xleftmargin=6pt, xrightmargin=6pt,
  breaklines=true, postbreak=\mbox{\textcolor{gris}{$\hookrightarrow$}\space},
  columns=fullflexible, keepspaces=true, showstringspaces=false,
  literate={á}{{\'a}}1 {é}{{\'e}}1 {í}{{\'i}}1 {ó}{{\'o}}1 {ú}{{\'u}}1
           {Á}{{\'A}}1 {É}{{\'E}}1 {Í}{{\'I}}1 {Ó}{{\'O}}1 {Ú}{{\'U}}1
           {ñ}{{\~n}}1 {Ñ}{{\~N}}1 {·}{{\textperiodcentered}}1
           {→}{{$\rightarrow$}}1 {…}{{\ldots}}1 {✔}{{\textcolor{verde}{$\checkmark$}}}1
}
\lstset{style=shell}

\newcommand{\cod}[1]{\texttt{\small #1}}
\DeclareUrlCommand\ruta{\def\UrlFont{\small\ttfamily}}
\setlength{\emergencystretch}{3em}

% Dónde se ejecuta cada bloque de comandos
\newcommand{\donde}[1]{\par\smallskip\noindent{\small\sffamily\bfseries\color{white}\colorbox{utaRojo}{\,#1\,}}\par\nopagebreak\vspace{-2pt}}
\newcommand{\equipo}{\donde{EN TU EQUIPO}}
\newcommand{\servidor}{\donde{EN EL SERVIDOR}}

\newtcolorbox{nota}[1][Nota]{colback=azulNota!5,colframe=azulNota,fonttitle=\bfseries,
  title=#1,boxrule=0.6pt,arc=2pt,left=6pt,right=6pt,top=4pt,bottom=4pt}
\newtcolorbox{aviso}[1][Importante]{colback=ambar!7,colframe=ambar,fonttitle=\bfseries,
  title=#1,boxrule=0.6pt,arc=2pt,left=6pt,right=6pt,top=4pt,bottom=4pt}
\newtcolorbox{esperado}{colback=verde!6,colframe=verde,boxrule=0.5pt,arc=2pt,
  left=6pt,right=6pt,top=3pt,bottom=3pt,fontupper=\small}

\renewcommand{\arraystretch}{1.25}
\newcolumntype{L}{>{\raggedright\arraybackslash}X}
\newcolumntype{P}[1]{>{\raggedright\arraybackslash}p{#1}}
\newcommand{\cabecera}{\rowcolor{utaRojo!12}}
\newcommand{\casilla}{$\square$\ }

\setcounter{tocdepth}{1}

\begin{document}

% ---------------------------------------------------------------------
\begin{titlepage}
  \centering
  \vspace*{2.5cm}
  {\color{utaRojo}\rule{\linewidth}{1.2pt}}\\[14pt]
  {\Huge\bfseries ServiceDesk UTA\par}
  \vspace{8pt}
  {\LARGE Tutorial de despliegue\par}
  \vspace{4pt}
  {\large Subir el sistema y ponerlo en marcha en otro servidor\par}
  \vspace{6pt}
  {\color{utaRojo}\rule{\linewidth}{1.2pt}}\\[30pt]
  \begin{tabular}{rl}
    \textbf{Versión del sistema:} & 1.0.0 (candidata a release) \\
    \textbf{Repositorio:} & \href{https://github.com/THThoms/servicedesk-fisei/tree/microservicios}{\ttfamily THThoms/servicedesk-fisei}, rama \texttt{microservicios} \\
    \textbf{Servidor destino:} & Ubuntu Server 24.04 LTS (x86\_64) \\
    \textbf{Fecha:} & 10 de noviembre de 2026 \\
  \end{tabular}
  \vfill
  {\small\color{gris}Para quien instalará el sistema en un servidor nuevo, sin necesidad de conocer el proyecto.\\
  Este documento no contiene contraseñas, secretos ni direcciones de servidores reales.\\
  Los valores entre \texttt{<...>} los sustituye cada instalación.\par}
\end{titlepage}

\pagenumbering{roman}
\tableofcontents
\clearpage
\pagenumbering{arabic}

% =====================================================================
\chapter{Antes de empezar}

\section{Qué vas a hacer}
El código está en GitHub. A partir de él se genera un \textbf{paquete de instalación}
(un archivo \cod{.tar.gz}). Ese paquete se copia al servidor, donde un instalador levanta
SQL Server, los tres microservicios, el gateway y el frontend en Docker.

\begin{center}
\begin{tikzpicture}[font=\small,
  caja/.style={draw=utaRojo,rounded corners=3pt,minimum width=3.1cm,minimum height=1.25cm,align=center,fill=white,line width=0.6pt},
  f/.style={-{Stealth[length=2.4mm]},line width=0.7pt,color=gris}]
  \node[caja] (git) {\textbf{GitHub}\\rama \texttt{microservicios}};
  \node[caja,right=2.1cm of git] (pkg) {\textbf{Paquete}\\\texttt{.tar.gz} + SHA-256};
  \node[caja,right=2.1cm of pkg] (srv) {\textbf{Servidor}\\Ubuntu + Docker};
  \draw[f] (git) -- node[above,font=\footnotesize]{1. generar} (pkg);
  \draw[f] (pkg) -- node[above,font=\footnotesize]{2. copiar} (srv);
  \node[below=0.35cm of srv,font=\footnotesize,align=center,text=gris] {3. instalar, 4. dominio y HTTPS,\\5. usuarios, 6. backups};
\end{tikzpicture}
\end{center}

Hay dos caminos para llevar el paquete al servidor. Los dos instalan exactamente lo mismo.

\begin{center}
\begin{tabularx}{\linewidth}{@{}P{4.3cm}LL@{}}
\toprule
\cabecera & \textbf{Camino A (recomendado)} & \textbf{Camino B} \\
\midrule
Dónde se genera el paquete & En tu equipo & En el propio servidor \\
Cómo llega al servidor & Se copia con \cod{scp} & Se clona el repositorio en el servidor \\
Ventaja & El servidor no necesita acceso a GitHub ni Git & No hace falta copiar archivos \\
Capítulo & \ref{cap:paquete}, sección \ref{sec:caminoA} & \ref{cap:paquete}, sección \ref{sec:caminoB} \\
\bottomrule
\end{tabularx}
\end{center}

\section{Qué necesitas}
\begin{center}
\begin{tabularx}{\linewidth}{@{}P{4.2cm}L@{}}
\toprule
\cabecera \textbf{Elemento} & \textbf{Detalle} \\
\midrule
Servidor & Ubuntu Server 24.04 LTS, x86\_64 (SQL Server no funciona en ARM). Mínimo 2 vCPU, 4\,GB de RAM y 40\,GB de disco para pruebas; para producción, 4 vCPU, 8\,GB y 100\,GB \\
Acceso & Usuario con \cod{sudo} y \textbf{acceso SSH con clave} (no con contraseña) \\
Dominio & Un nombre, por ejemplo \cod{servicedesk.tu-dominio.edu.ec}, y permiso para crear su registro DNS \\
Puertos & 80 y 443 accesibles desde Internet (Let's Encrypt los necesita), además del de SSH \\
Edición de SQL Server & \textbf{Pruebas}: Developer. \textbf{Producción}: Express (hasta 10\,GB por base) o una edición con licencia \\
Microsoft Entra ID & Quien administre el tenant debe añadir una \emph{Redirect URI} (capítulo \ref{cap:sso}) \\
Copia externa & Una cuenta de Backblaze B2 con un \emph{bucket} privado (capítulo \ref{cap:backups}) \\
Tu equipo & Git y Git Bash (en Windows), y un cliente SSH (\cod{ssh}, \cod{scp}) \\
\bottomrule
\end{tabularx}
\end{center}

\section{Cómo leer este tutorial}
Cada bloque de comandos indica dónde se ejecuta:

\equipo
\begin{lstlisting}
comando que se escribe en tu computadora (Git Bash en Windows)
\end{lstlisting}
\servidor
\begin{lstlisting}
comando que se escribe en el servidor, conectado por SSH
\end{lstlisting}

Debajo de los pasos importantes, un recuadro verde indica qué debe aparecer:
\begin{esperado}
Lo que tienes que ver si todo va bien. Si ves otra cosa, no sigas: consulta el capítulo \ref{cap:problemas}.
\end{esperado}

\begin{aviso}[Reglas de seguridad durante todo el proceso]
\begin{itemize}[leftmargin=*]
  \item Nunca subas a GitHub un archivo \cod{.env}, claves, \emph{keystores}, \cod{google-services.json} ni backups. El repositorio es \textbf{público}.
  \item No escribas contraseñas en documentos, chats ni capturas. El instalador las genera y las guarda en el servidor con permisos restringidos.
  \item No abras los puertos 1433, 8080 ni 3000 en el cortafuegos. Solo se usan dentro del servidor.
\end{itemize}
\end{aviso}

% =====================================================================
\chapter{Subir el código a GitHub}\label{cap:github}

El código consolidado ya está publicado en la rama \cod{microservicios} de
\url{https://github.com/THThoms/servicedesk-fisei/tree/microservicios}.
Este capítulo solo hace falta si has hecho cambios y quieres subirlos antes de desplegar.

\section{Revisar antes de subir}
\equipo
\begin{lstlisting}
cd <carpeta-del-proyecto>
git switch microservicios
git status                      # qué ha cambiado
git diff                        # revisa el contenido línea a línea
\end{lstlisting}

Comprueba que en la lista no aparece ningún \cod{.env}, \cod{seed.demo.ts}, \cod{*.keystore},
\cod{google-services.json}, \cod{*.bak} ni nada de \ruta{release/} o \ruta{backups/}. El
\cod{.gitignore} los excluye, pero conviene mirarlo.

\section{Confirmar y subir}
\equipo
\begin{lstlisting}
git add <archivos-que-cambiaste>
git commit -m "tipo(área): descripción breve del cambio"
git push origin microservicios
\end{lstlisting}

\begin{esperado}
\cod{microservicios -> microservicios} al final del \cod{push}, sin errores.
\end{esperado}

\begin{nota}[Si tu copia local no tiene el remoto configurado]
Añádelo una vez con\\
\cod{git remote add fisei https://github.com/THThoms/servicedesk-fisei.git}\\
y usa \cod{git push fisei microservicios}. Nunca uses \cod{git push --force} ni \cod{--mirror}.
\end{nota}

% =====================================================================
\chapter{Generar el paquete y llevarlo al servidor}\label{cap:paquete}

\section{Camino A: generar el paquete en tu equipo}\label{sec:caminoA}
El paquete se genera siempre desde un \textbf{clon nuevo}, no desde la carpeta en la que trabajas.
Así contiene exactamente lo que está en GitHub, con finales de línea de Linux (LF).

\equipo
\begin{lstlisting}
git clone -c core.autocrlf=false -b microservicios \
    https://github.com/THThoms/servicedesk-fisei.git sd-release
cd sd-release
bash scripts/build-server-package.sh 1.0.0
\end{lstlisting}

\begin{esperado}
Una serie de comprobaciones con \textcolor{verde}{$\checkmark$}, entre ellas:\\
\cod{sin auditorías, pruebas, cuentas demo, IP públicas ni cuentas reales} y\\
\cod{scripts y configuración del servidor con fin de línea LF}.\\
Al final:\\
\cod{release/servicedesk-uta-server-1.0.0.tar.gz (568K, 225 archivos)}\\
y su SHA-256. Si aparece una \textcolor{red}{\ding{55}}, el script se detiene y explica el motivo.
\end{esperado}

El paquete y su suma de comprobación quedan en \ruta{sd-release/release/}. Cópialos al servidor:

\equipo
\begin{lstlisting}
scp release/servicedesk-uta-server-1.0.0.tar.gz \
    release/servicedesk-uta-server-1.0.0.tar.gz.sha256 \
    <usuario>@<IP-o-nombre-del-servidor>:/tmp/
\end{lstlisting}

Después verifica en el servidor que el archivo llegó íntegro:

\servidor
\begin{lstlisting}
cd /tmp && sha256sum -c servicedesk-uta-server-1.0.0.tar.gz.sha256
\end{lstlisting}

\begin{esperado}
\cod{servicedesk-uta-server-1.0.0.tar.gz: OK}. Si dice \cod{FAILED}, vuelve a copiarlo.
\end{esperado}

\section{Camino B: clonar y generar en el servidor}\label{sec:caminoB}
\servidor
\begin{lstlisting}
sudo apt update && sudo apt -y install git
git clone -b microservicios \
    https://github.com/THThoms/servicedesk-fisei.git /tmp/sd-src
cd /tmp/sd-src
bash scripts/build-server-package.sh 1.0.0
cp release/servicedesk-uta-server-1.0.0.tar.gz* /tmp/
\end{lstlisting}

El resultado esperado es el mismo que en el camino A. A partir de aquí, los dos caminos continúan igual.

% =====================================================================
\chapter{Preparar el servidor}\label{cap:preparar}

\section{Ubuntu}
\servidor
\begin{lstlisting}
sudo apt update && sudo apt -y full-upgrade
sudo apt -y install ca-certificates curl gnupg rsync openssl nginx certbot ufw
sudo timedatectl set-timezone America/Guayaquil
\end{lstlisting}

Si el servidor solo trae el usuario \cod{root}, crea uno administrativo y copia tu clave SSH a ese usuario:
\begin{lstlisting}
adduser admin && usermod -aG sudo admin
\end{lstlisting}

\begin{aviso}[Antes de desactivar las contraseñas de SSH]
Comprueba en \textbf{otra} terminal que puedes entrar con tu clave. El instalador endurece
SSH (sin contraseña, sin \cod{root} por contraseña) solo si encuentra una clave autorizada, y lo valida
antes con \cod{sshd -t}. Aun así, no cierres la sesión actual hasta comprobar que puedes volver a entrar.
\end{aviso}

\section{Docker Engine y Compose}
Instálalos desde el repositorio oficial de Docker, no desde el paquete \cod{docker.io} de Ubuntu:

\servidor
\begin{lstlisting}
sudo install -m 0755 -d /etc/apt/keyrings
curl -fsSL https://download.docker.com/linux/ubuntu/gpg | \
    sudo gpg --dearmor -o /etc/apt/keyrings/docker.gpg
echo "deb [arch=$(dpkg --print-architecture) \
signed-by=/etc/apt/keyrings/docker.gpg] \
https://download.docker.com/linux/ubuntu \
$(. /etc/os-release && echo $VERSION_CODENAME) stable" | \
    sudo tee /etc/apt/sources.list.d/docker.list
sudo apt update
sudo apt -y install docker-ce docker-ce-cli containerd.io \
    docker-buildx-plugin docker-compose-plugin
sudo docker version --format '{{.Server.Version}}' && sudo docker compose version
\end{lstlisting}

\begin{esperado}
Una versión de Docker reciente (se recomienda 28 o superior) y \cod{Docker Compose version v2.x}.
\end{esperado}

El cortafuegos (UFW con SSH, 80 y 443) lo configura el instalador. Si tu SSH no usa el puerto 22,
indícalo luego con \cod{--ssh-port N}.

% =====================================================================
\chapter{Instalar ServiceDesk}\label{cap:instalar}

\section{Extraer el paquete}
\servidor
\begin{lstlisting}
sudo mkdir -p /opt/servicedesk
sudo tar -xzf /tmp/servicedesk-uta-server-1.0.0.tar.gz \
    -C /opt/servicedesk --strip-components=1
cd /opt/servicedesk && ls
\end{lstlisting}

\begin{esperado}
\cod{api-gateway auth-service catalog-service deploy docs frontend scripts ticket-service}, además de
\cod{MANIFEST.md}, \cod{FILES.txt} y \cod{VERSION}. Con \cod{ls -a} verás también \cod{.env.example}.
\end{esperado}

\section{Configurar \texttt{deploy/server/.env}}
\servidor
\begin{lstlisting}
cd /opt/servicedesk
sudo cp .env.example deploy/server/.env
sudo chmod 600 deploy/server/.env
sudo nano deploy/server/.env
\end{lstlisting}

Rellena estas variables:

\begin{center}
\begin{tabularx}{\linewidth}{@{}P{5.0cm}L@{}}
\toprule
\cabecera \textbf{Variable} & \textbf{Qué poner} \\
\midrule
\cod{PUBLIC\_DOMAIN} & El dominio, sin \cod{https://}. Ejemplo: \cod{servicedesk.tu-dominio.edu.ec} \\
\cod{SERVER\_PURPOSE} & \cod{pruebas} o \cod{produccion} \\
\cod{MSSQL\_PID} & \cod{Developer} solo si es de pruebas; \cod{Express} (o tu edición con licencia) en producción. Con \cod{produccion}, Developer se \textbf{bloquea} \\
\cod{INITIAL\_ADMIN\_EMAIL}, \cod{INITIAL\_ADMIN\_NAME} & Correo y nombre del primer administrador. \textbf{No} se pone contraseña \\
\cod{MICROSOFT\_*} (4 variables) & Datos del registro de la aplicación en Entra ID (capítulo \ref{cap:sso}). Las cuatro o ninguna \\
\cod{ADMIN\_EMAILS} & Correos que entrarán como administradores con Microsoft \\
\cod{EMAIL\_SEND\_MODE}, \cod{SMTP\_*} & \cod{log} para pruebas; \cod{smtp} con los datos del servidor de correo para enviar \\
\cod{FCM\_*} & Solo si se usarán notificaciones push en Android \\
\cod{BACKUP\_ALERT\_EMAIL} & Correo que recibirá avisos si falla un backup \\
\bottomrule
\end{tabularx}
\end{center}

\begin{aviso}[Deja vacías las contraseñas]
\cod{MSSQL\_SA\_PASSWORD}, las \cod{*\_DB\_PASSWORD} y \cod{JWT\_SECRET} deben quedar \textbf{vacías}:
el instalador las genera al azar con \cod{--generate-secrets} y no las muestra.
\end{aviso}

\section{Ensayo sin instalar nada}
\servidor
\begin{lstlisting}
sudo scripts/install-server.sh --mode clean --generate-secrets --dry-run
\end{lstlisting}

\begin{esperado}
Una tabla con cada variable como CONFIGURADA, VACÍA o AUSENTE (nunca su valor) y al final
\cod{Configuración válida}. Cada \ding{55} dice qué corregir en el \cod{.env}.
\end{esperado}

\section{Instalar}
Elige el modo:

\begin{itemize}
  \item \textbf{Instalación limpia}: sistema nuevo, sin datos.
\servidor
\begin{lstlisting}
sudo scripts/install-server.sh --mode clean --generate-secrets
\end{lstlisting}
  \item \textbf{Restauración}: trasladar los datos de otra instalación. Necesitas un backup completo
        (los tres \cod{.bak}, las imágenes y \cod{SHA256SUMS}) copiado en el servidor. Valídalo primero con \cod{--dry-run}:
\servidor
\begin{lstlisting}
sudo scripts/install-server.sh --mode restore --from /root/restore/<ID> \
    --generate-secrets --dry-run
sudo scripts/install-server.sh --mode restore --from /root/restore/<ID> \
    --generate-secrets
\end{lstlisting}
\end{itemize}

El instalador pide confirmar la licencia de SQL Server (escribe \cod{SI}) y, antes de tocarlos, el
cortafuegos y SSH. La primera vez tarda entre 5 y 15 minutos, porque compila las imágenes.

\begin{esperado}
Entre otras líneas: la edición real de SQL Server, \cod{6/6 servicios healthy} y
\cod{SMOKE: ... 0 fallos}. En modo restauración, además: \cod{SHA-256 correcto}, \cod{RESTORE VERIFYONLY},
\cod{DBCC CHECKDB} correctos y los conteos de usuarios, tickets e imágenes.
\end{esperado}

\section{Contraseña del administrador inicial}
Solo en la instalación limpia. El instalador guarda una contraseña \textbf{temporal} en un archivo
que únicamente puede leer \cod{root}. Léela una vez y bórrala:

\servidor
\begin{lstlisting}
sudo cat deploy/server/.initial-admin-secret
sudo shred -u deploy/server/.initial-admin-secret
\end{lstlisting}

Al entrar por primera vez, la aplicación obliga a cambiarla por una propia.

% =====================================================================
\chapter{Dominio y HTTPS}\label{cap:https}

\section{DNS}
En el proveedor del dominio, crea un registro \cod{A} con el nombre del dominio y la IP del servidor.
Comprueba que ya resuelve:

\servidor
\begin{lstlisting}
dig +short servicedesk.tu-dominio.edu.ec
\end{lstlisting}

\begin{esperado}
La IP del servidor nuevo. El cambio de DNS puede tardar unos minutos en propagarse.
\end{esperado}

\section{Certificado de Let's Encrypt}
\servidor
\begin{lstlisting}
cd /opt/servicedesk
sudo certbot certonly --webroot -w /var/www/certbot \
     -d servicedesk.tu-dominio.edu.ec \
     --agree-tos -m <correo-de-contacto> --no-eff-email
sudo scripts/install-server.sh --nginx-tls
sudo certbot renew --dry-run
\end{lstlisting}

\begin{esperado}
\cod{Successfully received certificate}, luego la configuración de nginx con TLS, y la prueba de
renovación correcta. La renovación es automática.
\end{esperado}

\section{Comprobar el sistema por el dominio}
\servidor
\begin{lstlisting}
sudo SMOKE_BASE_URL=https://servicedesk.tu-dominio.edu.ec \
     scripts/operations/smoke-server.sh
curl -s -o /dev/null -w '%{http_code}\n' http://servicedesk.tu-dominio.edu.ec/
\end{lstlisting}

\begin{esperado}
\cod{SMOKE: ... 0 fallos} y \cod{301}, la redirección de HTTP a HTTPS.
\end{esperado}

% =====================================================================
\chapter{Inicio de sesión con Microsoft}\label{cap:sso}

\begin{enumerate}
  \item Quien administra el tenant de Microsoft Entra ID añade esta \emph{Redirect URI} al registro
        de la aplicación:
        \begin{center}\cod{https://<PUBLIC\_DOMAIN>/api/auth/microsoft/callback}\end{center}
  \item En \ruta{deploy/server/.env} se ponen las cuatro variables \cod{MICROSOFT\_CLIENT\_ID},
        \cod{MICROSOFT\_TENANT\_ID}, \cod{MICROSOFT\_CLIENT\_SECRET} y \cod{MICROSOFT\_REDIRECT\_URI}.
        La última debe ser \textbf{exactamente} igual a la URL anterior.
  \item Se aplican los cambios:
\servidor
\begin{lstlisting}
cd /opt/servicedesk
sudo env RUN_SEED=false docker compose -p servicedesk \
     -f deploy/server/compose.yaml \
     --env-file deploy/server/.env up -d
\end{lstlisting}
  \item Se prueba abriendo \cod{https://<dominio>} y pulsando «Iniciar sesión con Microsoft».
\end{enumerate}

Solo se admiten cuentas del dominio de \cod{ALLOWED\_DOMAIN} (por defecto, \cod{uta.edu.ec}).

% =====================================================================
\chapter{Usuarios y roles}\label{cap:usuarios}

Todas las personas entran como \cod{user}. Los roles de técnico, DTIC y administrador se asignan desde el
servidor. Para asignar un rol a alguien que usa Microsoft, esa persona debe entrar antes al menos una vez.

\servidor
\begin{lstlisting}
cd /opt/servicedesk
# área Técnicos, área DTIC, bloque del técnico y comprobación:
sudo scripts/operations/manage-user.sh set-role <correo> technician
sudo scripts/operations/manage-user.sh set-role <correo> dtic
sudo scripts/operations/manage-user.sh set-block <correo> BLOQUE_1
sudo scripts/operations/manage-user.sh list-staff
\end{lstlisting}

\begin{center}
\begin{tabularx}{\linewidth}{@{}P{6.2cm}L@{}}
\toprule
\cabecera \textbf{Otras operaciones} & \textbf{Comando} \\
\midrule
Crear una cuenta solo con Microsoft & \cod{manage-user.sh create <correo> <rol> "<Nombre>"} \\
Crear una cuenta con contraseña local & Lo mismo con \cod{--password} (la pide sin mostrarla) \\
Recuperar el acceso de una cuenta & \cod{manage-user.sh reset-password <correo>} \\
\bottomrule
\end{tabularx}
\end{center}

\section{Prueba completa en el navegador}
\begin{enumerate}
  \item Como \textbf{usuario}: «Nuevo ticket», elige servicio y ubicación, escribe el detalle, adjunta una
        imagen JPG o PNG y créalo.
  \item Como \textbf{técnico} del área: «Disponibles», acepta el ticket y ábrelo desde «Mis aceptados».
        Debe verse la imagen.
  \item El técnico pulsa «Resolver». El usuario \textbf{rechaza} con un motivo y el ticket vuelve al técnico.
  \item El técnico resuelve de nuevo y el usuario \textbf{acepta}. El ticket queda cerrado con su constancia.
  \item Como \textbf{administrador}: revisa «Estadísticas», «Historial» y «Reportes», e imprime un reporte.
\end{enumerate}

% =====================================================================
\chapter{Copias de seguridad}\label{cap:backups}

El instalador deja programados el backup diario (03:15), su comprobación cada 6 horas, la retención
y la verificación semanales y un simulacro de restauración mensual. Falta conectar la
\textbf{copia externa cifrada} en Backblaze B2.

\section{Conectar Backblaze B2}
\begin{enumerate}
  \item En Backblaze, crea un \emph{bucket} \textbf{privado} y una \emph{Application Key} limitada a ese bucket.
  \item En el servidor, crea el archivo de credenciales, que está separado del \cod{.env}:
\servidor
\begin{lstlisting}
sudo install -d -m 700 /etc/servicedesk
sudo test -f /etc/servicedesk/backup.env || \
  sudo install -m 600 /opt/servicedesk/deploy/backup/backup.env.example \
    /etc/servicedesk/backup.env
sudo nano /etc/servicedesk/backup.env
\end{lstlisting}
        Rellena estas variables:
        \begin{itemize}
          \item \cod{RESTIC\_REPOSITORY}: \cod{s3:https://<endpoint>/<bucket>/servicedesk}
          \item \cod{AWS\_ACCESS\_KEY\_ID}: el \emph{keyID} de la Application Key
          \item \cod{AWS\_SECRET\_ACCESS\_KEY}: la \emph{applicationKey}
          \item \cod{B2\_REGION}: la región que aparece en el endpoint
        \end{itemize}
  \item Genera la clave de cifrado, inicializa el repositorio y comprueba la conexión:
\servidor
\begin{lstlisting}
cd /opt/servicedesk
sudo scripts/backup/restic-init.sh --generate-password
sudo scripts/backup/restic-init.sh
sudo scripts/backup/restic-test-connection.sh
\end{lstlisting}
\end{enumerate}

\begin{aviso}[Custodia de la clave]
La clave de cifrado queda en \ruta{/etc/servicedesk/restic.pass}. Sin ella, los backups externos
\textbf{no se pueden recuperar}. Guarda una copia fuera del servidor, en un gestor de contraseñas, y
nombra al menos dos personas que la custodien.
\end{aviso}

\section{Primer backup y simulacro}
\servidor
\begin{lstlisting}
sudo systemctl start servicedesk-backup.service
sudo scripts/operations/check-backup.sh
sudo scripts/backup/restore-drill.sh
\end{lstlisting}

\begin{esperado}
\cod{Backup completo}, \cod{Backup correcto de hace 0 h} y \cod{RESTORE DRILL PASS}. Si el backup
termina con código 3, falta la copia externa; con código 4, la copia externa falló.
\end{esperado}

% =====================================================================
\chapter{Comprobación final}\label{cap:final}

\servidor
\begin{lstlisting}
sudo scripts/operations/server-health.sh
\end{lstlisting}

\begin{esperado}
\cod{Resumen: 0 fallo(s)}. Los avisos se pueden revisar, pero no bloquean la puesta en marcha.
\end{esperado}

\section{Lista antes de abrir el sistema a los usuarios}
\begin{itemize}[label={}]
  \item \casilla Los 6 servicios aparecen como \cod{healthy} en \cod{docker compose ... ps}.
  \item \casilla La dirección \ruta{https://<dominio>/api/auth/config} muestra \cod{devLoginEnabled:false} y \cod{microsoftSimulateEnabled:false}.
  \item \casilla La contraseña temporal del administrador está cambiada y su archivo, borrado.
  \item \casilla \cod{sudo ufw status} solo permite SSH, 80 y 443.
  \item \casilla \cod{sudo ss -tlnp | grep -E ':(1433|8080|3000)'} solo muestra \cod{127.0.0.1}.
  \item \casilla SSH no acepta contraseñas: \cod{sudo sshd -T | grep passwordauthentication} da \cod{no}.
  \item \casilla En producción, la edición de SQL Server no es Developer y la licencia está confirmada.
  \item \casilla El certificado es válido y \cod{certbot renew --dry-run} es correcto.
  \item \casilla El backup manual terminó con copia externa y el simulacro de restauración pasó.
  \item \casilla El inicio de sesión con Microsoft funciona con una cuenta real.
  \item \casilla La prueba completa en el navegador fue correcta.
\end{itemize}

% =====================================================================
\chapter{Actualizar a una versión nueva}\label{cap:actualizar}

\begin{enumerate}
  \item \textbf{En tu equipo}: sube los cambios a GitHub (capítulo \ref{cap:github}) y genera el paquete nuevo
        con otra versión, por ejemplo \cod{1.0.1} (capítulo \ref{cap:paquete}). Cópialo y verifícalo en el servidor.
  \item \textbf{En el servidor}: haz un backup y aplica la versión nueva, conservando el \cod{.env}:
\servidor
\begin{lstlisting}
sudo systemctl start servicedesk-backup.service
sudo cp /opt/servicedesk/deploy/server/.env /root/servicedesk.env.respaldo
sudo tar -xzf /tmp/servicedesk-uta-server-1.0.1.tar.gz \
    -C /opt/servicedesk --strip-components=1
cd /opt/servicedesk
sudo env RUN_SEED=false docker compose -p servicedesk \
     -f deploy/server/compose.yaml \
     --env-file deploy/server/.env up -d --build
sudo scripts/operations/smoke-server.sh
\end{lstlisting}
  \item Las migraciones de la base de datos se aplican solas al arrancar cada servicio.
  \item Si algo falla, extrae de nuevo el paquete anterior y repite el último paso. Si la versión nueva
        había cambiado datos, restaura el backup del paso 2 con la guía de restauración y rollback
        que está en \ruta{docs/operations/}.
\end{enumerate}

\begin{nota}
El paquete no contiene el \cod{.env}, así que extraer una versión nueva no lo sobrescribe. Aun así,
la copia en \ruta{/root} protege contra un error al extraer.
\end{nota}

% =====================================================================
\chapter{Problemas frecuentes}\label{cap:problemas}

\begin{center}
\small
\begin{tabularx}{\linewidth}{@{}P{5.2cm}L@{}}
\toprule
\cabecera \textbf{Qué ves} & \textbf{Qué hacer} \\
\midrule
\cod{FAILED} al verificar el SHA-256 & El archivo se copió mal. Cópialo de nuevo \\
\cod{bash\textbackslash r: No such file or directory} & El paquete tiene finales de línea de Windows. Genéralo desde un clon nuevo (capítulo \ref{cap:paquete}) \\
El instalador rechaza \cod{MSSQL\_PID=Developer} & Estás en \cod{produccion}. Usa \cod{Express} o la edición con licencia \\
\cod{mssql} no llega a \emph{healthy} & Poca memoria (SQL Server necesita al menos 2\,GB). Revisa los logs con \cod{docker compose ... logs mssql} \\
Un servicio se reinicia una y otra vez & Revisa sus logs. Suele ser que las contraseñas del \cod{.env} se cambiaron después de instalar \\
Certbot no obtiene el certificado & El DNS aún no apunta al servidor o el puerto 80 está cerrado \\
Error de \emph{redirect URI} al entrar con Microsoft & La URI registrada en Entra ID no es idéntica a \cod{MICROSOFT\_REDIRECT\_URI} \\
Error de CORS en el navegador & \cod{PUBLIC\_DOMAIN} no coincide con el dominio desde el que se accede \\
Respuesta 429 al iniciar sesión & Demasiados intentos fallidos. Espera unos minutos \\
No se suben fotos del iPhone & El formato HEIC no se admite. Usa «Más compatible» en la cámara o comparte la foto como JPG \\
Backup con código 3 o 4 & Falta configurar B2 (código 3) o no hay conexión o las credenciales están mal (código 4) \\
\bottomrule
\end{tabularx}
\end{center}

La lista completa está en \ruta{docs/troubleshooting/PROBLEMAS_FRECUENTES.md}, dentro del paquete.

\begin{nota}[Estado de la validación]
Este procedimiento se ha validado instalando el paquete en una simulación de servidor nuevo: instalación
limpia, restauración, prueba funcional completa y recuperación ante caídas de SQL Server. Todavía no
se ha ejecutado en un Ubuntu real con systemd, nginx y certbot, ni con el inicio de sesión real de Microsoft,
correo SMTP o notificaciones push. La primera instalación real debe hacerse en un servidor de pruebas.
\end{nota}

\end{document}
